\section{Limitations}
\label{sec:limitations}

The screen is best interpreted as an investigative prioritization
device---a triage tool that flags suspicious procurement environments
for closer scrutiny---rather than evidence sufficient to infer
firm-level liability or sanction cartel membership. We document a
robust conditional association and a coherent diagnostic pattern; we
do not claim causal identification of the FL--price association or
proof of the cover-bidding mechanism. The main limitations are as
follows.

\paragraph{Diagnostics are not identification.}
The five supporting diagnostics are coherent with coordinated cover
bidding but do not uniquely identify the mechanism. Each admits
alternative readings in isolation, and the joint pattern, while
difficult to reconcile with a simple competitive account, does not
exclude all non-collusive explanations. The diagnostics strengthen
the case for deploying the screen; they do not constitute proof of
collusion, and they are not substitutes for forensic investigation.

\paragraph{Within-market selection.}
Item, year, and PBU fixed effects absorb time-invariant cross-market
heterogeneity, and the cross-fit removes any mechanical link between
classification and outcomes. The residual threat is time-varying
within-market confounders that attract both FL firms and higher
prices. The Cinelli--Hazlett bound ($RV = 17.5\%$) quantifies the
required confounder strength; the supporting diagnostics
(Section~\ref{sec:mechanisms}) constrain the set of stories
compatible with the data, but constraining is not identification.

\paragraph{Event-study pre-trends.}
Positive coefficients at $t = -2$ and $t = -3$ before first FL
entry are consistent with strategic market selection
(Section~\ref{sec:additional_id}) and equally compatible with
unobserved confounding. The pre-trends preclude a causal reading
of the DiD and reinforce the paper's non-causal framing; the paper
relies on within-item conditional comparisons for this reason.

\paragraph{Classification boundaries.}
The screen depends on a zero-win bright line and a participation
threshold. Relaxing either weakens the coefficient. Sophisticated
cartels could evade by allowing cover bidders occasional wins or
staying below the threshold---hence the three-stage workflow
(Section~\ref{sec:conclusion}) that combines FL screening with
bid-level forensics and periodic threshold recalibration. Any
operational deployment would require updating the threshold as
procurement environments evolve.

\paragraph{Screen scope.}
The FL screen flags suspicious environments, not individual firms.
It does not identify which firms within a flagged environment are
colluding, does not estimate the causal effect of collusion on
prices, and does not produce evidence sufficient for sanctions or
adjudication. These boundaries are inherent to the screen's design
as an investigative prioritization device rather than a forensic
tool.

\paragraph{External validity.}
All results come from BEC. The screen's minimal data requirements
make replication straightforward, but generalization to other
procurement platforms, institutional settings, and enforcement
regimes requires testing.
